\section*{第 \ref{ch:g9:earth-resources} 章——地球资源}

\begin{solution}{exo:g9:earth-resources:1}
(a) \omterm{def:g9:earth-resources:aggregate}{骨料}。(b) \omterm{def:g9:earth-resources:reserve}{储量}。(c) \omterm{def:g9:earth-resources:ore-deposit}{矿床}。(d) \omterm{def:g9:earth-resources:ore-deposit}{富集系数}。
\end{solution}

\begin{solution}{exo:g9:earth-resources:2}
(a) 位于\omterm{def:g1:landscapes-around-us:water}{河流}\omterm{def:g3:stones-pebbles-sand:pebble}{砾石}中。(b) 位于裂隙中，由热水形成。(c) 位于\omterm{def:g7:volcanoes:magma}{岩浆}中。
(d) 这是\omterm{def:g1:landscapes-around-us:water}{河流}\omterm{def:g3:stones-pebbles-sand:pebble}{砾石}\omterm{def:g9:earth-resources:ore-deposit}{矿床}形成的第一步：被释放出的金
随后被\omterm{def:g1:landscapes-around-us:water}{河流}搬运。
\end{solution}

\begin{solution}{exo:g9:earth-resources:3}
(b)、(d)、(a)、(e)、(c)。
\end{solution}

\begin{solution}{exo:g9:earth-resources:4}
(a) 不可再生：形成需要数百万年。(b) 可再生：
降雨补充被取走的部分。(c) 不可再生：\omterm{def:g9:earth-resources:ore-deposit}{矿床}的形成
需要极长的时间。(d) 按此速率开采则不可再生：\omterm{def:g6:water-underground:water-table}{地下水位}下降，
\omterm{def:g6:water-underground:spring}{水井}最终干涸。
\end{solution}

\begin{solution}{exo:g9:earth-resources:5}
(a) $4 \div 0.007 \approx 571$，约为 $10^3$（介于 $10^2$ 与 $10^3$ 之间，
更接近 $10^3$）。(b) $0.5 \div 0.0002 = 2500$，约为 $10^3$。(c) $20
\div 5.6 \approx 3.6$，约为 4：介于 $10^0$ 与 $10^1$ 之间。
(d) 铁在\omterm{def:g6:inside-the-earth:layers}{地壳}中已十分丰富：矿石的品位只需比普通\omterm{def:g3:stones-pebbles-sand:rock}{岩石}高几倍，
便已足够。
\end{solution}

\begin{solution}{exo:g9:earth-resources:6}
(a) $19.3 \div 2.5 \approx 7.7$ 倍。(b) 金：$19.3 \times 0.01 =
\qty{0.193}{g}$；\omterm{def:g3:stones-pebbles-sand:rock}{岩石}：$2.5 \times 0.01 = \qty{0.025}{g}$。(c) 整个岩层移动时，
较重的金颗粒会沉入较轻的沙和\omterm{def:g3:stones-pebbles-sand:pebble}{卵石}之间，
直至坚硬的\omterm{def:g3:stones-pebbles-sand:rock}{基岩}处才停下。
\end{solution}

\begin{solution}{exo:g9:earth-resources:7}
(a) 金、锡和铬。(b) 介于 $10^1$ 与 $10^2$ 之间，接近
$10^2$。(c) $3000 \div 3.7 \approx 800$ 倍。(d) 铬（或锡，或
金）：\omterm{def:g9:earth-resources:ore-deposit}{富集系数}越高，
大自然使该金属充分富集的地方就越罕见。
\end{solution}

\begin{solution}{exo:g9:earth-resources:8}
(a) $18 \times 365 = \qty{6570}{kg}$，约为 \qty{6.6}{t}。(b) $5 \times
10^{10}$ 吨。(c) $25000 \div 18 \approx 1389$ 天，接近四年。
(d) 从河中挖出的沙不再向下游输送到\omterm{def:g1:landscapes-around-us:water}{海岸}；
海浪持续将沙从海滩带走，而补充的沙却减少了，
以致无法替代。
\end{solution}

\begin{solution}{exo:g9:earth-resources:9}
(a) 来自\omterm{def:g2:water-on-earth:ocean}{海洋}中的微小动植物，如硅藻，它们埋藏在
海底的泥中。(b) 约 60 至 \qty{120}{\degreeCelsius} 之间。
(c) 它们比充满\omterm{def:g3:stones-pebbles-sand:rock}{岩石}\omterm{def:g6:water-underground:porosity}{孔隙}的水轻。
(d) 不透水层（如页岩）位于其上方；它们聚集在该层之下，常位于背斜的顶部，
气在上，油在下。
\end{solution}

\begin{solution}{exo:g9:earth-resources:10}
(a) $2.9 \times 10^7$ 吨。(b) $2.9 \times 10^7 \div 100 = 2.9 \times
10^5$ 吨每年，大致如此。(c) 矿石的质量是其所含铜的
100 倍：$2.9 \times 10^7 \times 100 = 2.9 \times 10^9$ 吨。
(d) 每吨矿石只能提取 \qty{10}{kg} 铜，因此
必须开挖数十亿吨\omterm{def:g3:stones-pebbles-sand:rock}{岩石}。
\end{solution}

\begin{solution}{exo:g9:earth-resources:11}
(a) 若每年开采 2300 万吨，该\omterm{def:g9:earth-resources:reserve}{储量}约可维持 $280 \div 23
\approx 12$ 年，即直至 20 世纪 80 年代初；然而实际开采量为 6.97 亿吨，
远超该\omterm{def:g9:earth-resources:reserve}{储量}。(b) $980 \div 280 = 3.5$ 倍。(c) $980 \div
23 \approx 43$ 年。(d) 例如：不断发现新\omterm{def:g9:earth-resources:ore-deposit}{矿床}；更好的技术或更高的价格
使品位较低的\omterm{def:g9:earth-resources:ore-deposit}{矿床}变得值得开采。
\end{solution}

\begin{solution}{exo:g9:earth-resources:12}
(a) $5 \times 10^1$ 年。(b) $2 \times 10^3 \div 50 = 40$ 倍。(c) $3
\times 10^8 \div (5 \times 10^1) = 6 \times 10^6$ 倍。(d) 煤形成所用的时间
是我们使用它所用时间的数百万倍：大自然无法
在我们的时间尺度上将其更新。
\end{solution}

\begin{solution}{pb:g9:earth-resources:1}
\textbf{1.} $100 \times 1000 = \num{100000}$~千克 $= 1 \times 10^5$~千克。

\textbf{2.} $2000 \times 300 = \num{600000}$~吨 $= 6 \times 10^5$~吨。

\textbf{3.} $6 \times 10^5 \div 10 = 6 \times 10^4$~吨每年。

\textbf{4.} $0.40 \times 6 \times 10^4 = 2.4 \times 10^4$~吨沙，
每年。

\textbf{5.} $1200 \times 500 \times 5 = 3 \times 10^6$~m$^3$。

\textbf{6.} 其中的三分之二：$2 \times 10^6$~m$^3$。

\textbf{7.} $2 \times 10^6 \times 1.6 = 3.2 \times 10^6$~吨。

\textbf{8.} $\num{24000} + \num{136000} = \num{160000}$~吨 $= 1.6 \times
10^5$~吨每年。

\textbf{9.} $\num{160000} \div 2000 = 80$ 倍。

\textbf{10.} 不能：采坑被掏空的速度比\omterm{def:g1:landscapes-around-us:water}{河流}补充\omterm{def:g3:stones-pebbles-sand:sand}{沙子}的速度快
80 倍。从城镇的尺度来看，这种\omterm{def:g3:stones-pebbles-sand:sand}{沙子}是有限资源。

\textbf{11.} $\num{24000} \div 4 = \num{6000}$~吨\omterm{def:g3:stones-pebbles-sand:sand}{沙子}，每年。

\textbf{12.} $3.2 \times 10^6 \div (1.6 \times 10^5) = 20$：\textbf{20 年}。
\end{solution}
